\chapter{Открытый банкинг, A2A и~ISO~20022}\label{ch:open-banking}
В~платеже без PAN объектами модели становятся банковский счёт, право доступа
к~нему, API для инициации операции и~согласие клиента. Здесь и~далее \textbf{A2A}
(\textit{account-to-account})~--- класс расчётной модели, в~которой деньги
переводятся напрямую между банковскими счетами, минуя карточные сети.

\section{PSD2 и~SCA}\label{sec:psd2-sca}

\importantnote{%
  \textbf{Применимость к~РФ.} Российское право использует иную модель:
  платёжные услуги регулирует 161-ФЗ, а~доступ к~счёту через открытые API
  развивается через стандарты Банка России и~АФТ. В~карточном канале функцию
  усиленной аутентификации выполняют 3-D~Secure и~механизмы противодействия
  мошенничеству эмитента (см.~гл.~\ref{ch:3ds}). Базовый акт надзора
  за~платёжными услугами~---
  Федеральный закон от~27.06.2011~№~161-ФЗ \enquote{О~национальной платёжной
  системе}~\cite{fz-161}
  (см.~гл.~\ref{ch:regulation}). Глава описывает PSD2 как систему, чтобы
  объяснить логику открытого банкинга и~ISO~20022; российская инфраструктура
требует отдельного анализа. }

\highlight{PSD2 (Payment Services Directive~2, Directive~(EU)~2015/2366)
  перевела доступ к~счёту из~частного банковского канала в~регулируемый
интерфейс}. Директива описывает платёж как денежное событие и~отдельно~--- роли,
через которые он инициируется. В~её словаре есть
\textbf{PSU (payment service user)}~--- клиент, дающий согласие; \textbf{ASPSP
(account servicing payment service provider)}~--- банк или другой провайдер,
ведущий счёт; \textbf{AISP} (\textit{account information service provider})~---
лицензируемая сторона для доступа к~данным счёта; \textbf{PISP} (\textit{payment
initiation service provider})~--- лицензируемая сторона для инициирования
платежа~\cite{psd2}. Совокупно AISP и~PISP называют \textbf{TPP}
(\textit{third-party provider})~--- сторонним поставщиком услуг по~PSD2.
Регулируемый внешний интерфейс становится допустимым способом начать операцию
со~счётом.

\textbf{Согласие} (\enquote{consent}) в~открытом банкинге~--- самостоятельный
объект: документированное волеизъявление PSU, фиксирующее объём, срок и~условия
доступа TPP к~данным счёта или к~инициации платежа. У~согласия собственный
жизненный цикл, отдельные статусы и~отдельный идентификатор, отличный
от~идентификатора платёжной операции.

Если платёж инициирован через PISP, директива требует от~него сразу после
инициации передать плательщику подтверждение и~ссылочный идентификатор операции.
Банку счёта PISP должен передать ссылочный идентификатор этой~же
операции~\cite{psd2}. Объект \enquote{инициация платежа} с~самого начала
получает собственный идентификатор и~сохраняет отдельный жизненный цикл.

SCA (\textit{Strong Customer Authentication}, усиленная аутентификация клиента)
добавляет к~этой схеме инициации платежа отдельный этап контроля. Банк
счёта убеждается, что плательщик подтверждает доступ к~счёту или платёж как
минимум двумя независимыми элементами из~трёх категорий: знание, владение,
биометрия~\cite{eba-rts-sca}. В~карточном канале смежную задачу решают
3-D~Secure и~собственные механизмы противодействия мошенничеству эмитента.
В~расчётах со~счёта SCA~--- базовое условие доступа.

Для удалённого электронного платежа RTS (\textit{Regulatory Technical Standards}
  EBA, Европейской банковской службы, по~SCA и~общим безопасным каналам связи~---
EBA~RTS on~SCA and~CSC) добавляют динамическое связывание (\textit{dynamic
linking}): до~подтверждения плательщик видит сумму и~получателя,
код аутентификации привязан именно к~этой сумме и~этому получателю, любое
изменение одного из~двух параметров аннулирует ранее сгенерированный
код~\cite{eba-rts-sca}. Смысл подтверждения~--- конкретная операция
с~конкретными денежными атрибутами.

SCA связывает между собой согласие клиента, банк счёта и~сам платёж. Экран
продукта и~экран банковского подтверждения должны показывать один набор
денежных атрибутов: счёт, получателя и~сумму. Расхождение этих атрибутов делает
результат SCA недостоверным для продукта.

RTS одновременно задают основания для упрощения SCA: малая сумма, доверенный
получатель, повторяющаяся операция на~одну и~ту же сумму одному и~тому же
получателю, перевод между собственными счетами у~одного провайдера и~анализ
транзакционного риска (TRA). Операции, инициированные продавцом, регулируются
отдельно от~этого перечня~\cite{eba-rts-sca}. Обязательное подтверждение
и~упрощение по~запросу~--- разные процедуры. \highlight{Решение по~исключению
принимает ASPSP}: третья сторона запрашивает упрощённый сценарий, а~банк счёта
выбирает итоговый режим.

PSD2 и~RTS задают распределение ответственности: банк счёта остаётся владельцем аутентификации,
согласия и~окончательного решения по~операции. В~карточной системе продавец
работает через эквайера и~PSP. В~открытом банкинге ASPSP становится явным
участником продуктовой архитектуры: он~управляет экраном согласия, проводит SCA
и~возвращает допуск или отказ; оба исхода предусмотрены протоколом.

Банки сопротивлялись PSD2 по~экономической причине: директива сокращала их
монополию на~клиентский интерфейс. До~PSD2 банк контролировал и~счёт,
и~интерфейс доступа к~нему. Каждый платёж через интернет-банк или мобильное
приложение генерировал данные, которые банк мог монетизировать (допродажи
(\textit{cross-sell}), кредитный скоринг, аналитика). PISP переносит инициацию
платежа в~свой интерфейс,
AISP~--- доступ к~данным счёта. Банк продолжает вести счёт и~нести расчётный
риск, а~контакт с~клиентом в~момент платежа и~эксклюзивность данных переходят
к~агрегатору. Добровольная модель API до~PSD2 ограничила массовый доступ
компаний финтеха; Еврокомиссия выбрала регуляторное принуждение
как способ создать конкурентный рынок~\cite{psd2}.

\section{Открытый банкинг: AISP, PISP и~поток согласия}\label{sec:open-banking}

Когда счёт становится доступным через регулируемый интерфейс,
\highlight{\textbf{AISP} и~\textbf{PISP} различаются в~доменной модели. AISP
  работает с~доступом к~данным счёта и~истории операций, тогда как PISP отвечает
за~инициацию конкретного платежа~\mbox{\cite{psd2}}}. В~карточном платеже PSP
объединяет эти функции в~продуктовой модели; в~открытом банкинге AISP и~PISP
хранятся и~проектируются раздельно.

Масштабируемой интеграции нужен общий профиль API для банков-участников.
В~Великобритании эту роль
выполняют стандарты Open Banking UK, в~континентальной Европе~--- Berlin Group
NextGenPSD2~\cite{openbanking-uk,berlin-group}. Open Banking UK публикует
профили Read/Write API, профиль безопасности и~рекомендации к~пользовательскому
опыту. Berlin Group описывает рамку XS2A с~моделями согласий, вариантами SCA
и~перечнем платёжных продуктов. От~выбора профиля зависит, где проверять
поведение банка~--- в~профиле API, профиле безопасности или документе
о~пользовательском потоке.

Типовой PISP-поток для одиночного платежа внутри страны в~Open Banking UK
(\textit{single domestic payment}~--- так этот класс операций называет
спецификация; рис.~\ref{fig:pisp-flow}) разделяет шаги по~двум фазам жизненного
цикла согласия на~платёж (\texttt{payment-consent} в~терминах API). На~фазе
создания PSU даёт явное согласие через ASPSP, выполняется SCA по~Article~97
PSD2~\cite{psd2}, согласие переходит в~статус
\texttt{Authorised}~\cite{openbanking-uk-payment-initiation-profile}. На~фазе
инициации PISP создаёт платёжный ресурс через \texttt{POST~/domestic-payments}
по~выданному \texttt{access\_token}, ASPSP исполняет перевод и~доводит статус
до~\texttt{AcceptedSettlement} или \texttt{Rejected}. Каждый одиночный платёж
требует собственного согласия с~собственной SCA. Серия последующих списаний
по~сохранённому мандату, аналогичная карточному \textit{card-on-file} MIT,
выносится в~другой класс операций.

Регулярные списания с~одной SCA в~Open Banking UK выделены в~отдельный
профиль Variable Recurring Payments (VRP)~\cite{openbanking-uk-vrp}. TPP создаёт
один объект VRP-согласия (\texttt{vrp-consent}) с~параметрами лимита (сумма,
частота, длительность), пользователь даёт SCA один раз, а~последующие платежи
VRP (\texttt{vrp-payment}) используют это согласие~--- через применение
Article~14 (\textit{recurring transactions}) или Article~13 (\textit{trusted
beneficiaries}) исключений SCA-RTS~\cite{eba-rts-sca}, либо через
\textit{delegated SCA} от~PISP.

\begin{figure}[htbp]
  \centering
  \includefiguresvg[width=\linewidth]{assets/figures/ch20-open-banking/pisp-flow}
  \caption{Типовой поток PISP в~OB~UK: создание согласия на~платёж
    с~SCA и~инициация одиночного платежа внутри
  страны.}\label{fig:pisp-flow}
\end{figure}
\FloatBarrier

В~этих API согласие и~платёж описаны как разные ресурсы. На~стороне AISP
типовой поток начинается с~создания ресурса согласия. В~Open Banking UK AISP
создаёт \texttt{account-access-consent}
с~использованием потока клиентских учётных данных
(\textit{client credentials grant}; поток OAuth~2.0, в~котором клиент
аутентифицируется перед сервером сам, без~участия пользователя), получает
\texttt{ConsentId} и~начальный статус \texttt{AwaitingAuthorisation}. После
аутентификации пользователя у~банка статус меняется на~\texttt{Authorised} или
\texttt{Rejected}; в~профиле версии v3.1.x отзыв согласия выделялся
как отдельное состояние \texttt{Revoked}. Для ресурсов, созданных начиная
с~v3.1.4, корректное конечное состояние~---
\texttt{Rejected}~\cite{openbanking-uk-account-access-consents,ob-uk-at-3110}.
Доступ к~данным счёта оформляется отдельным ресурсом согласия со~своим
жизненным циклом и~статусами.

У~PISP разделение видно на~статусах согласия и~платежа. В~типовом потоке Open
Banking UK PISP сначала создаёт ресурс \texttt{payment-consent} и~получает
\texttt{ConsentId}, после чего пользователь перенаправляется в~ASPSP для
аутентификации и~подтверждения, и~только после этого PISP создаёт платёжный ресурс.
Для ресурса согласия в~Open Banking UK заданы статусы \texttt{AwaitingAuthorisation},
\texttt{Rejected}, \texttt{Authorised} и~\texttt{Consumed}; \texttt{Consumed}
означает только использование согласия для создания платёжного ресурса.
Окончательное зачисление получателю подтверждает отдельный статус платежа
у~ASPSP~\cite{openbanking-uk-payment-initiation-profile,openbanking-uk-domestic-payment-consents}.
Финальность расчёта определяется отдельным статусом платежа. Команда, которая
ориентируется только на~экранное \enquote{платёж подтверждён}, смешивает эти
статусы.

Жизненный цикл \texttt{domestic-payment-consent}
(рис.~\ref{fig:consent-lifecycle}) уточняет правила перехода. После
\texttt{Authorised} переход \texttt{Revoke} закрыт (PSD2 Art.~80);
\texttt{DELETE} через ASPSP~--- отдельная операция до~\texttt{Authorised}
с~иным смыслом, чем \texttt{Revoked}~\cite{ob-uk-pi-3110}.

\begin{figure}[htbp]
  \centering
  \includefiguresvg[width=\linewidth]{assets/figures/ch20-open-banking/consent-lifecycle}
  \caption{Жизненный цикл \texttt{domestic-payment-consent} (PISP)
  в~Open Banking UK.}\label{fig:consent-lifecycle}
\end{figure}

Для аутентификации Open Banking UK поддерживает редирект (\textit{redirect})
и~развязанный подход. Редирект требует того~же устройства, того~же канала
и~полноценного браузера; телефония, POS и~устройства с~ограниченным экраном
требуют развязанного сценария. В~развязанном потоке SCA выполняет банк,
а~пользователь продолжает подтверждение в~другом клиенте
относительно исходного запроса~\cite{openbanking-uk-authentication-methods}.
Платёжному продукту требуется поддержка смены канала и~асинхронного возврата
результата; наличие редиректа задаёт только один вариант потока.

Поэтому продукт обязан явно хранить несколько объектов. Минимум включает
\texttt{consent\_id}, \texttt{payment\_id}, локальный идентификатор сессии
редиректа, срок жизни токена доступа, идентификатор операции у~банка, итоговый
статус платёжного ресурса и~момент окончания повторного использования согласия
или токена. При объединении этих идентификаторов и~состояний система присваивает
разным причинам один ярлык \enquote{обрыв возврата пользователя}: просроченный
\texttt{access\_token}, отклонённое согласие, отменённая у~банка операция или
отложенный финальный статус.

На~карточной странице оплаты (\textit{checkout}) продавец или его PSP управляет
взаимодействием до~ответа авторизации. В~открытом банкинге
пользователь подтверждает операцию в~интерфейсе банка, а~продукт получает результат
только через редирект, обратный вызов (\textit{callback}) или повторный запрос
статуса. Корректная модель содержит отдельный автомат состояний для согласия,
платёжного ресурса и~финального результата.

\section{Промышленные платформы открытого банкинга: ЕС и~США}\label{sec:open-banking-platforms-eu-us}

В~ЕС правила задаёт PSD2, в~США~--- рыночная инициатива, но~практическая
интеграция банковских счетов с~приложениями сформировала отдельный класс
продуктов~--- платформы-агрегаторы, которые лицензируются как TPP и~подключают
сотни банков под единым API. Их API используют большинство потребительских
сервисов финтеха.

Plaid~\cite{plaid-overview}~--- крупнейший агрегатор в~США: к~его API подключено
большинство известных потребительских приложений (Robinhood, Venmo, Coinbase,
Affirm). В~январе 2020~года Visa объявила о~приобретении Plaid за~\$5{,}3\,млрд;
в~ноябре того~же года Антимонопольное управление Минюста США подало иск
о~блокировке сделки, и~в~январе 2021~года стороны её
прекратили~\cite{plaid-visa-doj}. TrueLayer~\cite{truelayer-overview}~---
европейский PISP/AISP по~PSD2: авторизация FCA, паспортизация в~ЕС, продукты
Pay-by-Bank. Tink~\cite{tink-visa-acquisition} работает
на~европейском рынке агрегации открытого банкинга
(\textit{open banking aggregation}): в~марте 2022~года Visa завершила
покупку Tink за~1{,}8\,млрд евро (около~\$2{,}1\,млрд)~--- для Visa это была
вторая попытка войти в~сегмент после провала сделки с~Plaid.
Yodlee~\cite{yodlee-overview}~--- старейший игрок рынка, основан в~1999~году
в~США, в~2015~году куплен Envestnet за~\$660\,млн, а~в~2025~году продан частной
инвестиционной компании~STG.

Архитектурный контраст с~российской моделью: в~ЕС и~США существуют независимые
лицензируемые TPP, которые работают с~банковскими данными десятков и~сотен
банков-партнёров и~предъявляют продавцу единый API поверх фрагментированной сети
счетов; в~России Федеральный закон от~27.06.2011~№~161-ФЗ
\enquote{О~национальной платёжной системе}~\cite{fz-161} закрепляет другую
модель. Функцию платформы-агрегатора выполняют сами банки и~НСПК через
инфраструктуру СБП.

\section{Открытый банкинг в~России}\label{sec:open-banking-russia}

Российский открытый банкинг использует собственную модель внедрения.
В~ЕС доступ к~счёту задан директивой, обязательной для банков; в~России он
развивается как поэтапная инициатива Банка России и~сообщества финтеха
со~стандартизацией открытых API и~переходом от~рекомендации к~обязательству
для определённых сегментов рынка~\cite{cbr-open-api,afe-open-api}.

\textbf{Концепция открытых API.} 9~ноября 2022~года Банк России опубликовал
\enquote{Концепцию внедрения Открытых API на~финансовом рынке}: документ
описывает целевую модель со~стандартизированными API для обмена данными
и~инициирования операций со~счетами, едиными правилами информационной
безопасности и~поэтапным расширением круга участников
и~сценариев~\cite{cbr-open-api-concept}.

\textbf{Стандартизация и~роль АФТ.} Технические стандарты открытых API в~России
разрабатываются при участии Ассоциации развития финансовых технологий (АФТ)
\enquote{ФинТех}~--- профессионального сообщества, объединяющего банки,
компании финтеха и~Банк России. АФТ публикует рекомендованные стандарты OpenAPI
для разных сценариев: получение баланса и~выписки, инициирование платежа,
аутентификация клиента~\cite{afe-open-api}.\footnote{По состоянию на~декабрь
  2025~года Банк России опубликовал десять стандартов открытых API, разработанных
  на~площадке АФТ; они носят рекомендательный характер. Переход пилотных
  участников на~утверждённые версии спецификаций по~счетам физических
и~юридических лиц намечен на~1~октября 2026~года.}

\textbf{Связь с~СБП и~SBPay.} В~отличие от~европейской модели, где PSD2 и~SEPA
Instant~--- две отдельные инициативы с~разной логикой развития, российский
открытый банкинг изначально проектируется в~связке с~СБП. Платёжная инициация
через открытый API означает запуск перевода через СБП с~уже выработанным
согласием клиента; сервис SBPay (раздел~\ref{sec:sbp-sbpay})~--- ранний пример
продуктовой реализации этой модели для розничных сценариев.

\textbf{Различие моделей.} PSD2 в~ЕС~--- регуляторное обязательство банков
предоставить доступ к~счёту по~запросу TPP (с~авторизацией клиента).
В~российской модели внедрение обязательных API поэтапное:
сначала добровольное участие банков, потом обязательное~--- для определённых
типов API и~определённых участников. Право клиента отказать в~доступе
сохраняется в~обоих подходах, а~архитектура внедрения различается по~степени
регуляторного давления и~по~темпу.

\importantnote{%
  Европейская терминология AISP и~PISP применима к~российской модели
  как описание функций. Российская регуляторика закрепляет другую конструкцию:
  договор с~банком, стандартизированный механизм согласия и~участие
  в~информационном обмене по~Открытым API. Право инициировать платёж
  и~читать данные по~счёту в~России возникает из~договорных отношений
  участника с~банком и~стандартизированного механизма согласия.
}

\section{ФАПИ.СЕК: российский профиль безопасности открытых API}\label{sec:fapi-sec}

Концепция Банка России задаёт регуляторную модель; протокольный слой состоит
из~FAPI и~ФАПИ.СЕК. FAPI (Financial-grade API)~--- международный профиль
безопасности, на~котором построены открытый банкинг Великобритании и~часть
европейских внедрений. ФАПИ.СЕК~--- российская адаптация FAPI к~требованиям
информационной безопасности Банка России и~отечественной
криптографии~\cite{fapi-sec-2024}.

\subsection*{FAPI как профиль безопасности}

Базовый OAuth~2.0 (\textit{Open Authorization}~--- фреймворк делегированной
авторизации, RFC~6749) проектировался для авторизации доступа к~произвольным
API. Финансовый сценарий требует дополнительных гарантий: параметры запроса
передаются в~URL, ответ авторизационного эндпоинта возвращается неподписанным,
токен доступа отвязан от~ключа клиента и~TLS-сессии. В~фреймворке одни и~те же механизмы работают и~для входа
на~новостной сайт через Google, и~для авторизации платежа со~счёта клиента.
Финансовому рынку требуется более строгий профиль.

Профиль FAPI разрабатывает рабочая группа OpenID Foundation. Он расширяет
базовый OAuth~2.0 и~OpenID Connect (OIDC~--- уровень идентификации поверх
OAuth~2.0) до~уровня, пригодного для операций с~деньгами. \textbf{FAPI~1.0
Part~2: Advanced} (утверждён 12~марта 2021~года) обязал использовать подписанный
\textit{request object}~\cite{rfc-9101-jar} (JSON Web Token~--- JWT,
  RFC~7519~--- с~параметрами авторизационного запроса, подписанный закрытым ключом
  клиента; семейство стандартов JOSE: JWS~--- JSON Web Signature, JWE~--- JSON Web
Encryption, JWK~--- JSON Web Key) вместо параметров в~URL,
токены, ограниченные отправителем (\textit{sender-constrained}; токен доступа,
  используется только клиентом, доказавшим владение конкретным
криптоключом или TLS-сертификатом)
через mTLS (\textit{Mutual~TLS}~--- взаимная TLS-аутентификация, RFC~8705), JARM
(\textit{JWT Secured Authorization Response Mode}~--- подписанный JWT-ответ
авторизационного эндпоинта), аутентификацию клиента через
\texttt{tls\_client\_auth} или \texttt{private\_key\_jwt} и~PKCE (\textit{Proof
Key for Code Exchange}, RFC~7636~\cite{rfc-7636-pkce}) с~алгоритмом S256
(SHA-256-вариант code challenge). Профиль рассчитан на~конфиденциальных
клиентов~\cite{fapi-1-part2}.

\textbf{FAPI~2.0 Security Profile} (утверждён 19~февраля 2025~года) ужесточил
профиль. Pushed Authorization Requests (PAR, RFC~9126~\cite{rfc-9126-par})
стал обязательным: авторизационный запрос переносится из~URL на~отдельный
эндпоинт авторизационного сервера, который возвращает короткий
\texttt{request\_uri}. Токены, ограниченные отправителем, остались обязательными,
но~к~mTLS добавилась альтернатива в~виде DPoP (\textit{Demonstrating Proof of
  Possession}~--- привязка токена к~закрытому ключу клиента через подписанное
доказательство, RFC~9449~\cite{rfc-9449-dpop}) для сценариев со~сложной
инфраструктурой взаимного TLS. PKCE~S256 и~TLS~1.2+ остались
обязательными~\cite{fapi-2-security-profile}.

\textbf{FAPI~2.0 Message Signing} (утверждён 26~сентября 2025~года) дополнил
профиль обязательной подписью самих сообщений запроса и~ответа к~ресурсу;
подпись на~этапе авторизации покрывает только часть доказательной последовательности.
Подпись сообщений нужна для неотказуемости (\textit{non-repudiation}) в~сценариях, где
финансовый институт обязан доказать содержание операции спустя длительное время
после её исполнения~\cite{fapi-2-message-signing}.

Профили FAPI описывают, \emph{что} должны делать стороны и~\emph{как} выглядит
передаваемое сообщение. \emph{Чем} оно подписывается~--- RSA, ECDSA, ГОСТ,~---
задаёт реализующий профиль через алгоритмы JWS из~реестра IANA. Поэтому
национальная адаптация FAPI включает замену криптографической
части и~локальной регуляторной надстройке~--- логика протокола сохраняется.

\subsection*{ФАПИ.СЕК: что наследуется, что меняется}

СТО~БР ФАПИ.СЕК-1.6-2024, введённый приказом от~07.10.2024~№~ОД-1615
\enquote{О~введении ФАПИ.СЕК}, заменил редакцию 2020~года и~стал актуальным
профилем безопасности для открытых API. Стандарт разработан АФТ совместно
с~ИнфоТеКС при
участии Банка России~\cite{fapi-sec-2024,cbr-fapi-sec-9908}.

Структура стандарта прямо наследует профили FAPI. В~ФАПИ.СЕК выделены три
раздела: профиль безопасности OpenID API в~режиме только для чтения
(соответствует FAPI~1.0~Part~1: Baseline); профиль безопасности в~режиме чтения
и~записи (соответствует FAPI~1.0~Part~2: Advanced) и~JARM как защищённый
JWT-режим ответа авторизационного
эндпоинта~\cite{fapi-sec-2024}.\footnote{Международные спецификации, которые
наследует ФАПИ.СЕК:~\cite{fapi-1-part1},~\cite{fapi-1-part2},~\cite{oidf-jarm}.}
Расширенный профиль ФАПИ.СЕК требует тех~же механизмов, что и~FAPI Advanced:
подписанный \textit{request object}, взаимный TLS на~транспортном уровне,
токены, ограниченные отправителем, \texttt{private\_key\_jwt} либо
\texttt{tls\_client\_auth}, PKCE~S256.

Главное расхождение касается криптографии: ФАПИ.СЕК предписывает исключительное
применение отечественных криптографических алгоритмов согласно методическим
рекомендациям ТК~26 МР~26.2.002-2024 \enquote{Использование российских
  криптографических алгоритмов в~протоколах OpenID
Connect}~\cite{tc26-mr-26-2-002-2024}. Подпись JWT и~проверка подписи~---
по~ГОСТ~Р~34.10\nobreakdash-2012 на~эллиптических кривых вместо RSA или ECDSA. Хеш-функция
в~JWS, JWE и~JWT~--- ГОСТ~Р~34.11\nobreakdash-2012 \enquote{Стрибог} вместо SHA-2.
Симметричное шифрование в~JWE~--- по~ГОСТ~Р~34.12\nobreakdash-2015 (\enquote{Кузнечик} или
\enquote{Магма}) вместо AES. Транспорт~--- TLS с~российскими криптонаборами
(ГОСТ~TLS) вместо обычного TLS~1.2/1.3 с~международными шифросьютами. Подпись
запросов и~хранение ключей выполняются исключительно средствами
сертифицированных СКЗИ (средств криптографической защиты информации);
ФАПИ.СЕК допускает программные реализации криптографии только через
сертифицированное СКЗИ.

\begin{figure}[tbp]
  \centering
  \includefiguresvg[width=0.9\linewidth]{assets/figures/ch20-open-banking/fapi-sec-stack}
  \caption{Стек ФАПИ.СЕК: процедурная часть FAPI и~JARM поверх
  отечественной криптографии и~ГОСТ~TLS.}\label{fig:fapi-sec-stack}
\end{figure}

Отдельный стандарт СТО~БР ФАПИ.ПАОК-1.0-2024, введённый приказом
от~07.10.2024~№~ОД-1616 \enquote{О~введении ФАПИ.ПАОК}, описывает профиль
аутентификации, инициированной
клиентом по~отдельному каналу (Client Initiated Backchannel Authentication,
CIBA)~\cite{fapi-paok-2024}. Сценарий используется там, где у~устройства, инициирующего
доступ, ограничены браузер или экран подтверждения: оплата с~умной
колонки, подтверждение через мобильное банковское приложение в~ответ
на~инициацию из~контактного центра, аутентификация на~устройстве с~отдельным
клиентским входом. ПАОК относится к~ФАПИ.СЕК так~же, как развязанная аутентификация
(\textit{decoupled authentication})~--- к~основному сценарию с~редиректом
(\textit{redirect}) у~Open Banking
UK~\cite{openbanking-uk-authentication-methods}: разные каналы доставки шага
аутентификации при~общей модели доверия.

\subsection*{Согласие, аутентификация и~подпись запроса}

ФАПИ.СЕК наследует от~FAPI раздельную модель объектов. Согласие хранится
на~стороне ASPSP (в~российской
  терминологии~--- \textbf{ПУ}, \textit{поставщик услуг}; роль TPP исполняет
  \textbf{СПУ}~--- \textit{сторонний поставщик услуг}, обращающийся к~API~ПУ
от~имени клиента) с~собственными состояниями (создано, ожидает аутентификации
клиента, активно, отозвано, истекло). Сама аутентификация клиента выполняется
средствами банка согласно регуляторным требованиям. Подписанный запрос инициации
операции~--- отдельный объект, в~котором клиентское приложение доказывает, что
согласие действительно и~операция совпадает с~тем, что подтвердил клиент.

В~типовом сценарии инициации платежа (рис.~\ref{fig:fapi-sec-flow}) клиентское
приложение формирует \textit{request object}~--- JWT с~описанием операции, подписанный
закрытым ключом клиента средствами СКЗИ по~ГОСТ~Р~34.10\nobreakdash-2012. Приложение
отправляет объект на~эндпоинт PAR банка и~получает в~ответ короткий
\texttt{request\_uri}; этот обмен следует процедуре FAPI~2.0~Security Profile.
Браузер клиента перенаправляется в~банк по~этому \texttt{request\_uri}; банк
проверяет подпись запроса, выполняет многофакторную аутентификацию
(\textit{multi-factor}) по~требованиям к~средствам идентификации, показывает
экран согласия с~параметрами операции и~возвращает
\texttt{authorization\_code} по~обратному вызову.
Клиентское приложение меняет \texttt{authorization\_code}
на~\texttt{access\_token}, привязанный к~mTLS-сертификату или DPoP-ключу,
и~использует токен для самого вызова API инициации платежа.

Привязка \texttt{access\_token} к~mTLS-сертификату или DPoP-ключу ограничивает
токен контекстом ключа: для вызова API нужен сопутствующий криптоматериал,
который хранится у~легитимного клиента средствами СКЗИ. Повторное использование
токена с~другого узла инфраструктуры СПУ приводит к~отказу банка на~этапе вызова
API инициации платежа. Тот~же принцип защищает от~подмены подписи: запрос
инициации с~тем~же \texttt{access\_token} и~подписью JWT, отличающейся
от~открытого ключа клиента из~его документа Key~Set
(по~разделу~5.7.3.4 ФАПИ.СЕК~\cite{fapi-sec-2024}), банк отклоняет до~показа
экрана согласия как ошибку конфигурации клиента.

\begin{figure}[tbp]
  \centering
  \includefiguresvg[width=0.9\linewidth]{assets/figures/ch20-open-banking/fapi-sec-flow}
  \caption{Поток инициации платежа через ФАПИ.СЕК с~расчётом
  через ОПКЦ~СБП.}\label{fig:fapi-sec-flow}
\end{figure}

Каждый шаг процедуры даёт отдельный код ответа. Ошибка проверки подписи
\textit{request object} отклоняет запрос до~показа экрана согласия. Истекшее
или отозванное согласие отзывает \texttt{access\_token}; банк возвращает статус,
отличный от~\enquote{недостаточно средств} или \enquote{лимит превышен}.
Сертификат клиента в~mTLS, отличный от~привязки токена, запускает новый процесс
с~PAR. Продукт, который объединяет все эти отказы в~одну колонку
\enquote{общий сбой платежа}, ухудшает диагностику отказов так~же, как описано
в~разделе~\ref{sec:open-banking-boundaries} для Open Banking UK.

Жизненный цикл объекта согласия (рис.~\ref{fig:fapi-sec-consent-lifecycle})
зафиксирован в~разделе~9.3.3 \texttt{ConsentStatus} стандарта, введённого
приказом от~19.12.2025~№~ОД-2895 \enquote{О~введении стандарта согласия
на~доступ к~счетам юрлиц}~\cite{cbr-od-2895}. \texttt{DELETE}-операция доступна
как до~авторизации пользователем, так и~после~неё~--- это отличает российскую
модель от~Open Banking UK, где после \texttt{Authorised} \texttt{Revoke} закрыт.

\begin{figure}[htbp]
  \centering
  \includefiguresvg[width=0.8\linewidth]{assets/figures/ch20-open-banking/fapi-sec-consent-lifecycle}
  \caption{Жизненный цикл объекта согласия в~ФАПИ.СЕК.}\label{fig:fapi-sec-consent-lifecycle}
\end{figure}

Отношение ФАПИ.СЕК к~стандартам аутентификации в~карточном платеже отдельное.
SCA в~3-D~Secure аутентифицирует держателя карты
по~конкретной операции и~работает внутри карточной сети. ФАПИ.СЕК
аутентифицирует клиента банка по~доступу к~счёту и~работает поверх
инфраструктуры открытых API. У~двух механизмов общий принцип сильной
аутентификации из~двух факторов разных категорий и~разные точки доверия
в~операционном потоке. Для пользователя оба механизма сходятся к~уведомлению
в~приложении банка, а~регулируются и~сертифицируются раздельно.

\subsection*{Связка с~СБП и~график обязательности}

ФАПИ.СЕК~--- технический протокол доступа к~счёту и~инициации операций по~нему.
Когда инициация платежа доходит до~этапа исполнения, расчёт идёт через ту~же
инфраструктуру, что и~любой другой A2A-платёж в~России~--- через ОПКЦ~СБП
и~платёжную систему Банка России~(см.~раздел~\ref{sec:sbp-protocol}). Полная
A2A-модель состоит из~двух частей: регулируемого интерфейса доступа
к~счёту и~мгновенной межбанковской расчётной системы.

Концепция Банка России от~9~ноября 2022~года задала исходный
график~\cite{cbr-open-api-concept}, но~по~итогам обратной связи
от~166~организаций он был пересмотрен. Документ Банка России от~2~сентября
2024~года \enquote{Основные принципы и~этапы внедрения Открытых API
на~финансовом рынке} подтвердил гибридный подход и~перенёс начало обязательного
применения на~2026~год тремя этапами~\cite{cbr-open-api-stages-2024}.
До~начала обязательного применения в~2026~году использование Открытых API для
всех участников финансового рынка остаётся рекомендательным.

\textbf{Этап~1 (2026~год)} распространяет обязательные требования только
на~крупнейших участников трёх секторов финансового рынка~--- банковского,
страхового и~инвестиционного, по~критериям Банка России (все системно значимые
банки, крупнейшие брокеры и~страховые организации). Микрофинансовый сектор,
фигурировавший в~Концепции 2022~года, из~первого этапа
исключён~\cite{cbr-open-api-stages-2024}. Окончательный срок реализации этапа~1
зависит от~готовности правовых и~инфраструктурных условий; даты этапов~2 и~3
пересматриваются по~результатам этапа~1.

24~декабря 2025~года Банк России опубликовал обновлённый комплекс стандартов
Открытых API, утверждённых приказами от~19.12.2025.
Стандарты сохраняют рекомендательный характер; переход пилотных участников
на~новые версии спецификаций по~счетам физических и~юридических лиц запланирован
на~1~октября 2026~года~\cite{cbr-open-api-update-2025}. Окончательные сроки
и~состав обязательных API меняются; точные даты сверяются по~сайту Банка России
и~публикациям АФТ~\cite{cbr-open-api,afe-open-api,openbankingrussia-wiki}.

\subsection*{Целевой комплекс на~октябрь 2026}

Состав целевого комплекса зафиксирован приказами Банка России от~19.12.2025.
Приказ от~19.12.2025~№~ОД-2887 \enquote{О~введении общих положений
и~глоссария СТО~БР по~открытым API} утвердил общие положения
и~глоссарий~\cite{cbr-od-2887}. Приказ от~19.12.2025~№~ОД-2894
\enquote{О~введении стандарта доступа к~информации о~счетах физлиц} ввёл
Accounts API ФЛ версии~2.0~\cite{cbr-od-2894}. Приказы от~19.12.2025
№~ОД-2892 \enquote{О~введении стандарта согласия на~доступ к~счетам физлиц}
и~от~19.12.2025~№~ОД-2895 \enquote{О~введении стандарта согласия на~доступ
к~счетам юрлиц} ввели Account-Consents API~\cite{cbr-od-2892,cbr-od-2895}.
Параллельно в~комплекс входят правила взаимодействия (Accounts API Rules),
технический стандарт расширенного профиля безопасности и~рекомендательные
материалы~--- \textit{Best Practices} и~\textit{Methodological Recommendations}.
Все целевые стандарты вступают в~действие 1~октября
2026~года~\cite{openbankingrussia-wiki}.

В~российской модели европейские роли AISP и~PISP заменены собственной парой
ролей. \textbf{Поставщик услуг (ПУ)}~--- организация, ведущая счёт клиента
и~публикующая Открытые API; функциональный аналог ASPSP. \textbf{Сторонний
поставщик услуг (СПУ)}~--- организация, обращающаяся к~API ПУ от~имени клиента;
функциональный аналог TPP. И~ПУ, и~СПУ в~российской модели работают через
договор с~банком и~включение в~реестр участников информационного обмена
по~Открытым API, ведение которого планирует Банк России. Вместо европейского
лицензионного статуса действует договорная модель
с~реестром~\cite{cbr-open-api-stages-2024,openbankingrussia-wiki}.

Спецификация Account-Consents API задаёт три эндпоинта: \texttt{POST
/account-consents} для создания ресурса согласия, \texttt{GET
/account-consents/\{consentId\}} для проверки статуса и~\texttt{DELETE
/account-consents/\{consentId\}} для отзыва. Авторизация всех трёх вызовов~---
по~потоку \textit{Client Credentials}; обязательные заголовки~---
\texttt{x-fapi-interaction-id} (корреляция операций по~UUID)
и~\texttt{x-jws-signature} (отделяемая подпись JWS тела запроса
по~ГОСТ)~\cite{cbr-od-2895}. \begingroup\sloppy\emergencystretch=5em
\texttt{permissions} согласия наследуют модель Open
Banking UK: \texttt{ReadAccountsBasic}, \texttt{ReadAccountsDetail},
\texttt{ReadBalances}, \texttt{ReadTransactionsBasic},
\texttt{ReadTransactionsCredits}, \texttt{ReadTransactionsDebits},
\texttt{ReadTransactionsDetail}~\cite{cbr-od-2895}.\par\endgroup

Инициирование платежей через Открытые API выделено в~отдельную разрабатываемую
спецификацию (Payment Initiation API). Она относится к~следующим этапам
и~сохраняет рекомендательный статус~\cite{openbankingrussia-wiki}. До~её
введения единственный массовый канал A2A-инициации платежа со~счёта
в~России~--- СБП.

Управление согласиями клиентов в~среде Открытых API консолидируется в~Платформе
коммерческих согласий (ПКС), которую Банк России и~Минцифры строят на~базе
Единого портала государственных и~муниципальных услуг (Госуслуги). Хранение
самих клиентских данных остаётся у~участников рынка, на~ПКС~--- только согласия.
Пилот ПКС был запланирован на~2025~год с~последующим переходом к~централизованному
управлению жизненным циклом
согласий~\cite{cbr-open-api-stages-2024,cbr-open-api}.

\subsection*{Российский профиль рядом с~европейскими}

Сравнение ФАПИ.СЕК с~Open Banking UK, Berlin Group NextGenPSD2 и~базовым
FAPI~2.0 показывает, где российский профиль отличается от~европейских.

\begin{table}[tbp]
  \caption{ФАПИ.СЕК рядом с~европейскими профилями открытого
  банкинга.}\label{tab:fapi-sec-comparison}
  \begingroup
  \tablesetup
  \setlength{\tabcolsep}{3pt}
  \begin{tabularx}{\textwidth}{
      @{}B{0.18\textwidth}
      Y
      Y
      Y
      Y@{}
    }
    \toprule
    \headrow
    \thd{Параметр} & \thd{ФАПИ.СЕК (РФ)} & \thd{Open Banking UK} & \thd{Berlin Group (ЕС)} & \thd{FAPI~2.0 (OIDF)} \\
    \midrule
    Базовый профиль & FAPI~1.0 Advanced + ГОСТ-надстройка & FAPI~1.0 Advanced & собственный профиль на~FAPI и~OAuth~2.0 & FAPI~2.0 \\
    \rowsep
    Криптография & ГОСТ~Р~34.10/\allowbreak 34.11/\allowbreak 34.12 в~JWS, JWE, JWK & RSA, ECDSA, SHA-2 & RSA, ECDSA, SHA-2 & задаётся реализующим профилем \\
    \rowsep
    Транспорт & ГОСТ~TLS & TLS~1.2+ & TLS~1.2+ & TLS~1.2+ \\
    \rowsep
    Подпись запроса & request object на~ГОСТ, СКЗИ обязательно & signed JWT, mTLS & signed JWT, mTLS либо QSeal & PAR обязательно \\
    \rowsep
    Лицензируемые роли & договор с~банком + соответствие стандарту & AISP, PISP по~PSD2 & AISP, PISP, CBPII по~PSD2 & за~пределами профиля \\
    \rowsep
    Расчётный канал & ОПКЦ~СБП & Faster Payments & SEPA Instant либо локальная RTGS & за~пределами профиля \\
    \rowsep
    Статус & рекомендательный, поэтапный переход к~обязательному & обязательный для девяти крупнейших банков & обязательный по~PSD2 & добровольная адаптация \\
    \bottomrule
  \end{tabularx}\par
  \endgroup
\end{table}

Главное отличие ФАПИ.СЕК от~Open Banking UK и~Berlin Group инфраструктурное:
процедуры наследуются от~FAPI, а~под~ними
меняется криптография и~TLS на~отечественные алгоритмы, требуются
сертифицированные СКЗИ у~всех участников, лицензируемые роли TPP вытеснены
договорной моделью с~банком. Для команды, проектирующей продукт на~ФАПИ.СЕК,
FAPI~1.0~Advanced и~JARM остаются основным источником понимания процедурной
части, а~МР~26.2.002-2024 ТК~26, Положение от~17.08.2023~№~821-П
\enquote{О~защите информации при~переводах денежных средств}~\cite{cbr-821p}
и~Положение от~30.01.2025~№~851-П \enquote{О~противодействии переводам
без~согласия клиента}~\cite{cbr-851p}
(см.~раздел~\ref{sec:regulation-cbr-provisions})~--- источниками требований
к~криптографии и~эксплуатации.

Состав обязательных API, тарифы пилотов, точные сроки перехода и~актуальная
редакция приказов меняются и~сверяются по~сайту Банка России и~публикациям АФТ.
Книга фиксирует архитектуру; оперативную справку ведут Банк России и~АФТ.

\section{A2A-платежи}\label{sec:a2a-payments}

PISP инициирует перевод со~счёта клиента, а~расчётная система исполняет движение
денег между банками. Открытый банкинг даёт согласие и~безопасный доступ
к~счёту; A2A задаёт модель движения денег между банковскими счетами.

Сравнение национальных систем мгновенных платежей (СБП, Pix, UPI, FedNow, SEPA
Instant)~--- в~главе~\ref{ch:instant-payments-world}. A2A здесь~--- класс
расчётной модели; архитектура у~каждой страны своя.

В~Pix SPI~--- центральная инфраструктура расчёта мгновенных платежей
Banco Central do Brasil, валовой расчёт в~реальном времени через специальные
счета участников в~центральном банке, после расчёта перевод становится
безотзывным~\cite{pix-spi}. Положительный статус A2A-перевода значит, что деньги
уже зачислены получателю, а~карточная авторизация подтверждает только разрешение
на~будущее списание; подробный разбор финальности, возвратов и~споров~---
в~разделе~\ref{sec:sbp-vs-cards}, расхождение жизненного цикла двух моделей
сведено в~табл.~\ref{tab:card-vs-a2a}.

\begin{table}[H]
  \caption{Жизненный цикл карточной операции и~A2A-перевода.}\label{tab:card-vs-a2a}
  \begingroup
  \tablesetup
  \begin{tabularx}{\textwidth}{
      @{}B{0.18\textwidth}
      Y
      Y
      Y@{}
    }
    \toprule
    \headrow
    \thd{Система} & \thd{Что даёт первый положительный ответ} &
    \thd{Когда деньги считаются окончательными} &
    \thd{Как устроены возврат и~спор} \\
    \midrule
    Карточный платёж &
    Разрешение на~авторизацию и~резервирование средств &
    После отдельного клиринга и~расчёта по~правилам сети &
    Возврат существует отдельно от~авторизации/\allowbreak подтверждения; возвратный платёж встроен
    в~сетевые правила \\
    \rowsep
    A2A-платёж &
    Подтверждение инициирования перевода или его исполнения,
    в~зависимости от~системы &
    Когда межбанковский перевод исполнен и~система считает его окончательным &
    Возврат~--- новый перевод по~инициативе плательщика; спор решается
    правилами конкретной системы \\
    \bottomrule
  \end{tabularx}\par
  \endgroup
\end{table}

В~электронной коммерции мгновенная расчётная система сокращает для продавца
интервал между подтверждением покупателя и~зачислением денег. Промежуточный
статус требует отдельного правила отгрузки, а~возврат~--- отдельного денежного
перевода или процедуры спора. A2A требует своей модели управления заказом,
своей поддержки и~собственной логики сверки.

\section{ISO~20022 в~A2A и~переводах}\label{sec:iso20022}

ISO~20022 задаёт семантическую модель: у~поля есть имя и~место в~общей
структуре сообщения. Стандарт удобен там, где один и~тот же платёж должен
одновременно поддерживать маршрутизацию, санкционный контроль и~ПОД/ФТ, сверку,
выписку и~отчётность~\cite{iso-20022}.

Семейства сообщений соответствуют разным участкам платёжной последовательности. \texttt{pain}
покрывает клиентскую инициацию (\textit{customer initiation})~--- запросы
клиента на~инициирование платежа; \texttt{pacs} используется в~межбанковском
клиринге и~расчёте (\textit{payments clearing and settlement}); \texttt{camt}
относится к~управлению денежными средствами (\textit{cash management})
и~отвечает за~статусы, выписки, уведомления и~связанные сообщения управления
ликвидностью~\cite{iso-20022-business-areas}.

В~практической последовательности клиент или его система формирует \texttt{pain.001} как
инструкцию на~кредитовый перевод. На~межбанковском уровне перевод идёт как
\texttt{pacs.008}; ответ о~состоянии на~этом же уровне приходит через
\texttt{pacs.002}. Когда деньги отражаются в~учёте клиента, банк присылает
\texttt{camt.054} или включает операцию в~выписку вида \texttt{camt.053}.
У~каждого шага своя аудитория и~свой авторитетный статус. Если приложение
получает только пользовательский API, а~сверка строится по~\texttt{camt}, между ними
требуется явное внутреннее соответствие.

Структурированное сообщение окупается на~сверке. Для сверки и~расследования
нужны сумма, валюта, поля уровня \texttt{EndToEndId}, идентификаторы дебитора
и~кредитора, структурированное назначение платежа и~код причины в~статусном
сообщении. В~карточном канале эти данные распределяются между внешним
процессинговым диалектом (\textit{processing}) и~файлами клиринга. В~ISO~20022 они
закладываются в~саму модель сообщения.

Стоимость миграции включает больше, чем переход с~одной XML-схемы на~другую.
При переходе на~новый профиль ISO~20022 меняются правила заполнения полей,
правила применения (\textit{usage guidelines}), допустимые статусы
и~обязательность отдельных атрибутов. Каноническая модель платежа внутри
продукта локализует такой переход; общая внутренняя запись заставляет
одновременно переписывать API, оркестрацию и~сверку.

В~SWIFT период сосуществования двух форматов сообщений
(\textit{coexistence period}) для межбанковских платёжных инструкций~--- MT
(\textit{Message Type}, формат FIN~MT в~SWIFT) и~ISO~20022~--- завершился
в~выходные переключения (\textit{cutover weekend}) 21--23~ноября
2025~года~\cite{swift-iso20022-milestone,swift-iso20022-implementation}.
После этой даты основным рабочим форматом для таких инструкций стал ISO~20022.
С~1~января 2026~года резервная
обработка для отправителей, остающихся на~MT (\textit{contingency processing}),
и~сквозная конверсия MT$\leftrightarrow$ISO внутри сети~SWIFT
(\textit{in-flow translation}) тарифицируются отдельно. MT~199 и~MT~299 для
обработки исключений (\textit{exception handling}) сохраняются после ноября
2025~года как
переходный канал и~подлежат окончательной отмене в~ноябре 2027~года, когда их
функцию полностью переймут сообщения camt.110
и~camt.111~\cite{swift-iso20022-implementation}. Помимо названия сообщения
интегратору приходится учитывать профиль, правила переходного периода
и~процедуру обработки исключений.

В~России Банк России ведёт отдельный раздел по~стандарту ISO~20022, описывает
его как основу собственных стандартов финансовых сообщений и~отдельно
поддерживает обмен сообщениями ISO~20022
в~СПФС~\cite{cbr-iso20022-album,cbr-spfs-iso20022}. Структурированные
данные нужны для SWIFT, международных переводов и~внутренних платёжных систем,
которые используют счёт вместо PAN.

\FloatBarrier
\section{Открытый банкинг: отличия от~A2A и~карточной оплаты}\label{sec:open-banking-boundaries}

\highlight{В~разговорах об~A2A-платежах смешивают разные вещи:
  доступ к~счёту и~управление согласием; расчётный канал, по~которому идут
  деньги; язык межсистемного сообщения. Это разделение делает разговор
об~\mbox{\enquote{оплате со~счёта}} точным}.

\begin{table}[tbp]
  \caption{Три составляющие A2A-платежа: доступ к~счёту, расчёт,
  сообщение.}\label{tab:account-based-layers}
  \begingroup
  \tablesetup
  \setlength{\tabcolsep}{3pt}
  \begin{tabularx}{\textwidth}{
      @{}B{0.18\textwidth}
      Y
      Y
      Y@{}
    }
    \toprule
    \headrow
    \thd{Компонент} & \thd{На какой вопрос отвечает} & \thd{Кто управляет} &
    \thd{Ошибка при~смешении} \\
    \midrule
    API доступа и~согласия &
    Кто и~на каких правах может обратиться к~счёту &
    ASPSP, профиль API, профиль безопасности, управление согласиями &
    Команда принимает истечение токена или отказ согласия за~сбой
    расчётной системы \\
    \rowsep
    Расчётная система &
    Как деньги переводятся между счетами и~когда они окончательны &
    Платёжная система, расчётная инфраструктура, правила системы &
    Промежуточный статус принимается за~расчёт; возврат проектируется
    как отмена уже исполненного перевода \\
    \rowsep
    Стандарт сообщения &
    Как операция описывается между системами и~в~отчётности &
    профиль ISO~20022, правила применения, локальные альбомы сообщений &
    Внешний API и~внутренний реестр обязательств становятся зависимыми от~одного
    конкретного формата сообщения \\
    \bottomrule
  \end{tabularx}\par
  \endgroup
\end{table}

Карточная страница оплаты работает с~PAN, эмитентным каналом авторизации,
клирингом и~правилами споров карточной сети. A2A-модель работает с~согласием
на~доступ к~счёту, инициированием банковского перевода и~финальностью расчётной
системы. СБП~--- частный случай A2A-системы; настоящая глава вводит более общий
класс, в~котором доступ к~счёту и~банковский перевод образуют самостоятельную
архитектуру.

Последнее разделение~--- между отказом, статусом и~спором. Отказ на~этапе
согласия или инициации~--- предусмотренная часть протокола: токен истёк, клиент
отменил согласие, банк отклонил выбранный сценарий SCA, счёт недоступен. Статус
платежа~--- отдельная задача: деньги остаются промежуточными даже при
\texttt{Consumed} у~согласия или \texttt{Accepted} у~API. Спор после исполнения
перевода~--- ещё одна задача со~своей процедурой разрешения конфликта. Если
продукт смешивает эти задачи в~одну колонку \enquote{общий сбой платежа}, он
ухудшает диагностику и~поддержку.

Три инженерных объекта~--- согласие, денежная операция, межсистемное
сообщение~--- сохраняют собственные жизненные циклы и~после того, как
пользователь увидел экран \enquote{Оплачено}. Если эти состояния объединены
в~одну запись во~внутреннем учёте, истёкший токен, ожидание \texttt{pacs.002}
и~спор после расчёта попадают в~одну колонку ошибок и~требуют ручного разбора.

\practicenote{%
  При проектировании A2A-платежей разделите внутренние объекты: объект согласия,
  объект денежной операции и~объект межсистемного сообщения. Тогда миграция
  на~новый банк, новую расчётную систему или новый профиль ISO~20022 меняет
один объект; остальные части системы сохраняют свои контракты. }
